\chapter{Corrections to earlier interpretations}\label{app:a}

Many readings in the project record were later corrected, and the chapters give the current account. The tables below
list earlier readings that still appear elsewhere in the repository, each with where it was made and where it was
corrected, so that a stale statement can be recognized. Where an old reading is likely to be met, the chapter also
carries a short "Correction" note.

\section{The device, registers, encoding and metadata (chapters 1, 2, 3, 11)}\label{sec:device-registers-encoding-2}

\begin{xltabular}{\linewidth}{@{}>{\hsize=1.29\hsize}X>{\hsize=0.77\hsize}X>{\hsize=1.26\hsize}X>{\hsize=0.69\hsize}X@{}}
\caption{Corrected readings about the device, registers, encoding and metadata.}\label{tab:corr-device-encoding}\\
\toprule
Old reading & Where & Current reading & Where \\
\midrule
\endfirsthead
\tablecontinued{4}
\toprule
Old reading & Where & Current reading & Where \\
\midrule
\endhead
R126 and R127 read zero and drop writes & recon section 134 (MM~\mm{19}) & Instructions naming R126–R143 are squashed whole; a zero-filled buffer
hid it & MM~\mm{10}, MM~\mm{19}, MM~\mm{25.18} \\
A released register reads zero & MM~\mm{0.15}, MM~\mm{17} rule 37 & 0 in every controlled run; not established as a rule; read only where 0
and the old value agree & MM~\mm{25.95}, MM~\mm{25.145}, MM~\mm{25.147} \\
Releases take effect at the end of the instruction (as a general rule) & MM~\mm{25.145} & True of measured ALU forms; the word store reads its value after its
index's release & MM~\mm{25.117} \\
\op{798} holes at registers 2 and 3 mod 8 & digest (not MM) & R18, R19, R26, R27, R50, R51, R58, R59 low halves & MM~\mm{25.196} \\
\op{1004} holes are d = 14, 15 & MM~\mm{25.136.6} & R14, R15, R30, R31, R46, R47, R62+ & MM~\mm{25.133}, \code{agxforge/g17/cc.py} \\
The register file is 32 (ALU destination five bits; memory fields four
bits) & early \code{agxforge/g17/cc.py} field map & R0–R125 measured; fields re-measured at seven and eight bits & \code{agxforge/g17/cc.py} (hireg sweep), MM~\mm{10} \\
IRGPR32 is "a file no instruction writes" & MM~\mm{25.72} / row T4 & its users write ordinary Rn; the class is R0–R143 plus IR0-IR15 & MM~\mm{25.59}, Apple's class table \\
16 bytes need both byte10 bits 7 and 0 (class 7) & \code{docs/agx-isa-status.md} & True for its population; \code{\_WIDE7\_EXT} rule 12 + 4 × byte10{[}0{]}
for byte6 in \{0x60, 0x68, 0x70, 0x78\}; neither general & MM~\mm{25}, corpus spans \\
"Our decode agrees with Apple's framing 400/400" as
independent evidence & MM~\mm{25.63} & That decode calls Apple's decoder & MM~\mm{25.209} \\
D1 975, D2 690; only \op{5106}/\op{5107}
admitted & \code{docs/g17-isa-cartography.md} & D1 7,549 (\val{d1} with the 26A434 decoder), D2 \val{d2}; ten MMA opcodes admitted & \code{docs/g17-isa-reference.md}, MM~\mm{2} \\
Store "consumer token" with a wait counter & recon sections 15, 22–24 (MM~\mm{19}) & bits 5:4 of the source register number & MM~\mm{19} \\
MMA high bits are a result ring & recon section 14 (MM~\mm{19}) & eight-bit wait mask over load slots & MM~\mm{6}, MM~\mm{19} \\
Tensor load slot is operand 1 bits 20–27 & MM~\mm{6} (and chapter 9.1's table) & bits 20–23 slot + 1; bits 24–31 the load's wait mask & \code{agxforge/g17/memenc.py}, decoder flips \\
Slot 31 and the variable record describe register spill & MM~\mm{25.113} & They describe the constant program & MM~\mm{25.120} \\
The constant program's uniforms: "6 + number of writes",
"1 + highest \op{592} uniform" & MM~\mm{25.120} rounds 2, 3b & L3: 1 + max(K // 2) over every writer & MM~\mm{25.120} \\
76 inert class constants (list itemising 92 bytes) & ledger \code{g17-class-constants-are-mostly-inert} & 76 inert; the descriptor tail is not one of the 148 & MM~\mm{25.148} \\
LC\_TABLE (0x31) & project tooling & LC\_NOTE & MM~\mm{25.148} \\
\code{manifest.register\_count} is the program's & delivered manifests & 222 of 247 carry a template's value; read the container & MM~\mm{25.147} \\
The encoder may accept bytes whose field arithmetic reaches zero
residual & campaign encoder & must decode back; it once rounded an odd register silently & MM~\mm{2}, MM~\mm{23.3} \\
\bottomrule
\end{xltabular}

\section{Scalar arithmetic, control flow, special registers and cross-lane
operations (chapters 4, 5, 6, 8)}\label{sec:scalar-arithmetic-control-2}

\begin{xltabular}{\linewidth}{@{}>{\hsize=1.31\hsize}X>{\hsize=0.75\hsize}X>{\hsize=1.21\hsize}X>{\hsize=0.72\hsize}X@{}}
\caption{Corrected readings about scalar arithmetic, control flow, special registers and cross-lane operations.}\label{tab:corr-scalar-control}\\
\toprule
Old reading & Section & Current reading & Section \\
\midrule
\endfirsthead
\tablecontinued{4}
\toprule
Old reading & Section & Current reading & Section \\
\midrule
\endhead
\op{11375} field code 1 is a ==
b (cc \code{CSEL\_CC['eq']}) & MM~\mm{25.90} "Was" column & (a \& b) != 0, a bit test; printed code 15 & MM~\mm{25.90}; \code{ledger/claims.toml} \\
\op{1062} is fsat & MM~\mm{25.90} "Was" column & \opn{1062} is cross-lane dfdx.fine.sat; fp32 saturate is \opn{904} & MM~\mm{25.90} \\
\op{11179}/10 operand 2 inherited as 4
(unsigned) for every int-to-float & MM~\mm{25.90} "Was" column & operand 2 is signedness, 4 unsigned, 5 signed & MM~\mm{25.90} \\
A read\_sr's first consumer must wait (read\_sr is a
late value) & MM~\mm{6}, MM~\mm{0.8}, MM~\mm{25.90} & slot-0 reads are readable at distance 0; Apple's wait is
a free policy & MM~\mm{25.117}, MM~\mm{25.121} \\
A store cannot take a special register as its source (the value "has not
arrived") & ledger 2026–09-08, MM~\mm{25.117} first claim & the failure was the aliased store release of \op{17229} operand 6 & MM~\mm{25.117} \\
0x82 = quadgroup\_index\_in\_threadgroup, 0x98 =
simdgroups\_per\_threadgroup & \code{isa/g17-scalar-isa.toml} & 0x82 = \code{SR\_SIMD\_ELEM} (lane), 0x98 = \code{SR\_TG\_X\_SIZE}; the builtins are
computed from them & MM~\mm{0.8}, \code{agxforge/g17/mdgen.py} \\
\code{SR\_SIMD\_GRP} has no generalized slot-29 entry & MM~\mm{0.8} & entry 53 & MM~\mm{25.93}, \code{agxforge/g17/mdgen.py} \\
A released register reads zero & MM~\mm{0.15}, MM~\mm{17} rule 37 & zero in every controlled measurement; checker assumes zero or old value & MM~\mm{25.95}, MM~\mm{25.145} \\
The branch is a 4-byte word with a 12-bit displacement & \code{isa/g17-scalar-isa.toml} & 10-byte branch; the 4-byte reading is its first four bytes & \code{agxforge/g17/cf.py}, \code{isa/g17-exec-mask.txt} \\
Codes 11 and 15 mean not-equal & \code{isa/g17-condition-codes.toml} (prediction) & 15 is the bit-test pattern; 11 unseen & same file (measured row) \\
(uint)x lowers as a 34-instruction decomposition (gave 5 for −5.5) & MM~\mm{25.186} & one \op{9320}, mode 1 (gives 0) & MM~\mm{25.186}, MM~\mm{25.196} \\
\op{14060} might zero-extend into the
high half & MM~\mm{25.145} & it writes only its half & MM~\mm{25.145.4} \\
\bottomrule
\end{xltabular}

\section{Memory and synchronization (chapters 7, 9)}\label{sec:memory-synchronization-chapters-2}

\begin{xltabular}{\linewidth}{@{}>{\hsize=1.26\hsize}X>{\hsize=0.81\hsize}X>{\hsize=1.24\hsize}X>{\hsize=0.69\hsize}X@{}}
\caption{Corrected readings about memory and synchronization.}\label{tab:corr-memory-sync}\\
\toprule
Old reading & Where & Current reading & Where \\
\midrule
\endfirsthead
\tablecontinued{4}
\toprule
Old reading & Where & Current reading & Where \\
\midrule
\endhead
The device load carries no displacement field; address = base + (index +
displacement) × sizeof & \code{isa/g17-addressing-rules.json} & Displacement field exists, in bytes, unscaled; the element-count add
applies to unsigned \code{tid + k}; the (index + displacement) × element
law is the immediate-address family's & \code{isa/g17-address-formula.txt}, MM~\mm{25.141.18}, MM~\mm{25.177} \\
Imageblock store operand 1 bit 31 selects shared / coordinate-addressed
storage & MM~\mm{25.89} (round 9), MM~\mm{25.104.3} rounds 8–9, MM~\mm{25.145.2},
\code{tools/g17emu.py} comments & Bits 24–31 are the wait mask; bit 31 waits on slot 7 & MM~\mm{25.96}, MM~\mm{25.104} \\
Cross-threadgroup hand-off requires a dispatch boundary & MM~\mm{0.14} & Fence-count-fence works inside one dispatch; waiting patterns need
residency & MM~\mm{25.140.5}, MM~\mm{25.141.1} \\
A threadgroup barrier with mem\_device orders device memory across
threadgroups & cc's \code{barrier("device")} use, MM~\mm{25.140.5} & It does not; \op{14156} on both sides & MM~\mm{25.140.5} \\
\code{435/10} is "array", \code{447/6} is "threadgroup atomic" & MM~\mm{25.75} & \op{435}; \op{447} & glossary, MM~\mm{25.192} \\
\op{10094}'s
return fills slot 7 & glossary & Fill slot is operand 1 bits 20–23; Apple uses 0/1, this compiler 7 & MM~\mm{25.115} \\
Simultaneous uniform atomics combine and return a pre-combine value & MM~\mm{25.115} (struck) & Consumer waited on the wrong slot & MM~\mm{25.115} \\
\op{10094} operand 6 is a source
lifetime ("16 writes nothing") & ledger \code{g17-the-uniform-atomic-needed-one-lifetime-bit} & Operand 6 is the byte offset & MM~\mm{25.140.3} \\
\op{10090} has eight unassigned raw
operation codes & \code{isa/g17-atomic-semantics.json} & Retracted; codes depend on the form's field placement & \code{isa/g17-atomic-family.toml} \\
\op{621} = (0xFFFFFFFF
>{}> src) \& 0xF, saturating at 32 & \code{isa/g17-address-formula.txt} (first fit) & mask of min(max(W - src, 0), 4) ones, W = operand 5 & same file, resolved section \\
Vector displacement is in units of eight bytes & \code{asm.encode\_vec4} & Bytes, for loads & MM~\mm{25.141.18} \\
The 8-byte \op{17229} has no wait bit & \code{agxforge/g17/cc.py} & Apple's 8-byte \opn{17229} carries operand-1 wait bits;
this compiler's encoder does not write them & \code{isa/g17-wait-mask-census.json}, MM~\mm{25.146} \\
Only the accelerator's own reads have faulted & MM~\mm{17} rule 34 & Through Metal; below Metal, launch-state errors fault & \code{docs/g17-first-dispatch-trials.md} \\
A "three-bit store limit" on wait masks & an earlier recon section, recalled in MM~\mm{6} & Encoder-table artifact; \op{17256} uses all eight
bits & MM~\mm{6} \\
Load-to-use about 121~ns & MM~\mm{25.97} & L1 at idle clock; 30.5~ns warm (numbers in ch 12) & MM~\mm{25.121} \\
\bottomrule
\end{xltabular}

\section{The tensor unit and performance (chapters 10, 12)}\label{sec:tensor-unit-performance-2}

\begin{xltabular}{\linewidth}{@{}>{\hsize=1.32\hsize}X>{\hsize=0.75\hsize}X>{\hsize=1.32\hsize}X>{\hsize=0.61\hsize}X@{}}
\caption{Corrected readings about the tensor unit and performance.}\label{tab:corr-tensor-performance}\\
\toprule
Old reading & Section & Current reading & Section \\
\midrule
\endfirsthead
\tablecontinued{4}
\toprule
Old reading & Section & Current reading & Section \\
\midrule
\endhead
An fp32 value above 464 gives the (positive) canonical NaN in e4m3fn & MM~\mm{0.11} & An fp8 overflow keeps its sign (0xff / 0xfc for negatives); only NaN
inputs become positive canonical NaN & MM~\mm{25.105}, MM~\mm{25.105.1} \\
Hot loops: about 26 KiB, roughly 2,800 instructions & MM~\mm{0.1}, MM~\mm{7} & The step is at a 16 KiB loop body (15,971 fast); no cost at 26–64
simdgroups per core & MM~\mm{19}, MM~\mm{25.144.12} \\
Instruction-footprint cliff at (16,066, 17,158{]} whole-kernel bytes & recon 156 via MM~\mm{19} & The variable is the loop body, bracketing 16 KiB & MM~\mm{19} \\
Hardware counters: timestamps only (read as "no counters exist") & MM~\mm{0.15}, MM~\mm{25.48}, MM~\mm{25.141.6} & True of Metal's public API;
Instruments' limiter template records 89 counters & MM~\mm{25.144.4} \\
Limiter counters not reachable without Apple's GUI or
entitlement ("M4 status") & MM~\mm{25.144.4} (earlier bullet) & Reachable via \code{xctrace} with the GUI-saved template & MM~\mm{25.144.4} (later recording) \\
Power was not measured & MM~\mm{20} & GPU-rail power measured on both engines & MM~\mm{25.56}, MM~\mm{25.60} \\
H5 (destination keep flag) open & MM~\mm{20} & Set implies the next MMA writes the same tuple; hardware-inert & MM~\mm{6}, MM~\mm{25.22} \\
Saturation with transposes, chains, high registers not run & MM~\mm{20} & Saturation holds under transpose; clips once per issue in a chain & MM~\mm{25.42}, MM~\mm{25.89} \\
fp32 operand subnormals flush to zero & recon 5 via MM~\mm{19}; \code{tensorops\_model.py} line 12 docstring & Low 13 bits cleared only; subnormals kept & MM~\mm{5}, MM~\mm{19} \\
\op{5100} / \op{5101} never compiler-emitted & recon 27, 104 via MM~\mm{19} & Emitted under \code{relaxed\_precision} & MM~\mm{19}, MM~\mm{17} rule 29 \\
Destination layout is a property of the opcode & recon 100 via MM~\mm{19} & One map; \code{pos\_b(k,c) = pos(rotl1(k), c)}; the difference is the
producing kernel's row labeling & MM~\mm{3}, MM~\mm{19} \\
The hardware performs the matrix-to-fragment transposition & MM~\mm{19} & No tile-layout unit; ALU address arithmetic makes the layout & MM~\mm{7}, MM~\mm{19} \\
MMA high flag bits are an eight-slot result ring & recon 14 via MM~\mm{19} & An 8-bit wait mask over eight load slots & MM~\mm{6}, MM~\mm{19} \\
Two MMA hazards (stored result to register consumer; back-to-back MMAs) & recon 112 via MM~\mm{19} & Host dispatch-scoping (\code{dim}) artifact & MM~\mm{19}, MM~\mm{17} rule 5 \\
\code{matmul2d} allows one \code{run()} per kernel; chain must cross a
kernel boundary & recon 108, 121 via MM~\mm{19} & Slice offsets misread; single-kernel chain works & MM~\mm{19} \\
Feed relabelings are instruction properties & MM~\mm{0.10}, MM~\mm{13}, MM~\mm{17} rules 9–11 (unscoped text) & Only for kernels that pack A and D canonically; under this
compiler's B-row packing the feeds read the logical operand (\cref{sec:coordinate-systems}) & MM~\mm{25.103} \\
Two reduction axes follow different orders & recon 120 via MM~\mm{19} & One rule: per-lane chain then butterfly over the axis lane bits & MM~\mm{12}, MM~\mm{19} \\
Reduction sums run descending & recon 19 via MM~\mm{19} & Direction belongs to the compiled object (fast-math) & MM~\mm{12}, MM~\mm{19} \\
Legacy MMA and tensor unit share one issue slot & recon 31 via MM~\mm{19} & Complete overlap at saturation; partial overlap is in-order issue & MM~\mm{9}, MM~\mm{19} \\
Legacy MMA overlaps fp32 fma 81–86 percent & MM~\mm{19} (first version of recon 152) & True of \op{838} only; class-118 forms
0.37 or less & MM~\mm{9}, MM~\mm{19} \\
Occupancy is slot-limited, not register-limited (20–116 registers) & recon 167 via MM~\mm{19} & Direct counter: 44.6 → 32.1 simdgroups per core over
22–126 registers (Apple-compiled kernel) & MM~\mm{19}, MM~\mm{25.49} \\
Legacy bimodality: two-state 9/8 rate; 79.4~ns quantum & recon 158, 164 via MM~\mm{19} & Per-core staircase; bimodality only in boundary windows & MM~\mm{25.122} \\
Tensor-load element code 4 "means 32-bit" withdrawn & MM~\mm{25.5/R7} "R7 narrowed" & Element field = width in bytes; 4 = fp32 & MM~\mm{25.37} \\
\op{612}: second immediate is an absolute
bound & MM~\mm{19} & It is a width & MM~\mm{7}, MM~\mm{19} \\
\op{435} and \op{437} are
shifts & (glossary history) & 16-bit ANDs & MM~\mm{25.28} / glossary \\
Software-pipelining gain ≥ 1.2× preregistered & MM~\mm{25.97} & 1.08× at 8 simdgroups, 0.99× at 12.8 per core & MM~\mm{25.97}, MM~\mm{25.107} \\
Apple M512 58.6~µs; ratios 1.17×, 2.1–4.5× & MM~\mm{25.98} & 106.6~µs; 1.06–1.07× at equal geometry / M512 & MM~\mm{25.107} \\
WS/TP crossover about 64 tokens & MM~\mm{16} (model) & fp8 48, bf16 64 on the retained harness & MM~\mm{25.126} \\
L1 load latency about 121~ns & MM~\mm{25.97} & 30.5~ns at a warm clock & MM~\mm{25.121} \\
Residency used-125: \textasciitilde43.6 per core, then
\textasciitilde58; S = 1,024 shortfall & MM~\mm{25.115} (v3-v5) & Warm: 1,024/1,024; this compiler's counter 64–68 per
core at 125 & MM~\mm{25.115}, MM~\mm{25.121} \\
\code{ac\_bench} per-op figures in milliseconds & MM~\mm{25.55} & Seconds & MM~\mm{25.56} \\
Prefill GEMMs at 60–69 percent of peak & MM~\mm{25.155} & 80–89 percent in graph with a closure-checked profile & MM~\mm{25.160} \\
Constant-drain fusion model (6.54~µs per pipeline change) & MM~\mm{25.142.4} (preregistration) & Refuted; removing a change saved 0.6–2.5~µs & MM~\mm{25.142.4} \\
mlx-lm baseline via bench loop (1.18–1.26×) & MM~\mm{25.153}-25.181 & \code{generate\_step} baseline: 1.14–1.22× & MM~\mm{25.190} \\
\bottomrule
\end{xltabular}

\section{Superseded counts and wordings}\label{sec:accounting-moved-chapters}

The chapters no longer carry the reconciliation detail collected here: counts that changed between measurements,
inventories that differ between sources, and wordings that were later replaced. Each row gives where the earlier
statement was made and where it was settled.

\subsection*{The device, registers, encoding and metadata}\label{sub:device-registers-encoding}

\begin{xltabular}{\linewidth}{@{}>{\hsize=1.32\hsize}X>{\hsize=0.67\hsize}X>{\hsize=1.24\hsize}X>{\hsize=0.77\hsize}X@{}}
\caption{Superseded counts and wordings about the device, registers, encoding and metadata.}\label{tab:acct-device-encoding}\\
\toprule
Old reading or accounting & Where & Current reading & Where \\
\midrule
\endfirsthead
\tablecontinued{4}
\toprule
Old reading or accounting & Where & Current reading & Where \\
\midrule
\endhead
"Everything was measured on macOS 26.6.2 (25G83)" / "every figure \dots{}
one OS build" & \cref{sec:part-software}; earlier 
\cref{sec:measured-part-software} & 25G83 is the stated platform and is recorded by the below-Metal records;
the evidence files of the Metal-path measurements do not record an OS
build & \cref{sec:measured-part-software}, MM~\mm{25.91} \\
"A spinning threadgroup is never preempted"; "a grid barrier among
resident threadgroups is sound" & earlier \cref{sec:forward-progress-in-order}; MM~\mm{25.141.1}
wording & observations of one probe family in the tested configurations;
restrictions: do not depend on preemption; size waiting grids to
measured residency & \cref{sec:forward-progress-in-order}, MM~\mm{25.141.1} \\
A released register "reads zero" (MM~\mm{0.15}), "every later read
returns zero" (MM~\mm{6}, MM~\mm{17} rule 37), every release observed
read exactly 0 (MM~\mm{25.95}), against "zero or its old value, observed
both ways" (MM~\mm{25.145}), "observed only sometimes" (MM
\mm{25.145.4}), "returns 0 - sometimes, not always"
(\code{agxforge/g17/releasecheck.py}) & MM~\mm{0.15}, MM~\mm{6}, MM~\mm{17}, MM~\mm{25.95}, MM~\mm{25.145}, MM~\mm{25.145.4} & every controlled probe read zero; readback not established in general;
read a released register only where zero and the old value give the same
answer & \cref{sec:release-does} \\
\code{IR*} is "a distinct 16-register file"
(\code{agxforge/g17/registerdomain.py}) against the 160-member decoder
class IRGPR32 & \code{agxforge/g17/registerdomain.py}, MM~\mm{25.59} & the two agree. IR32 is the 16-register IR file; IRGPR32 is R0–R143
followed by IR0-IR15, and corpus users write only its general members & \cref{sec:register-classes} \\
\op{798} holes are "registers 2 and 3 mod 8"
(R34, R35, R42, R43 are not holes); \op{1004}
holes only d = 14 and 15 ("r119/r120" in decoder ids) & MM~\mm{25.136.6} & the listed sets: the fma's eight low-half destinations
and two sources; the widening's R14, R15, R30, R31, R46,
R47, R62 and above & \cref{sec:encoding-holes}, MM~\mm{25.133}, MM~\mm{25.196} \\
Instructions are "2 to 16 bytes" & \cref{sec:has-been-recovered}; the 216-opcode census
(\code{isa/g17-opcodes.toml}) & 2 to 22 bytes, even: Apple's corpora show every even
length 2–16 plus 18 and 22 (\op{10909} at 18, \op{14661} at 22), no 20; the decoder-admitted universe shows
4–16, 18, 22 & \cref{sec:framing} \\
"Apple's table declares 17,779" (MM~\mm{24.8} row T1,
\code{docs/g17-isa-cartography.md}, \code{tools/g17admit.py}) against
"17,796" (\code{docs/agx3-oracle.md}, \code{agxforge/g17/model.py}) & those sources & both hold for the 25G83 build, which has 17,796 table entries, 17,779 declaring at least one operand,
17 declaring none (\code{isa/g17-agx3meta-instrs.txt}) & \cref{sec:opcode-identity-counts} \\
D2 690, then 1,107 when per-form evidence was generated, then 1,118
through hardware checks of drifting bits & MM~\mm{25.87}, MM~\mm{25.58} & D2 = \val{d2} of D1 7,549 (\val{d1} with the 26A434 decoder) & \cref{sec:opcode-identity-counts} \\
D1 975, D2 690, and only \op{5106} and \op{5107}
admitted & \code{docs/g17-isa-cartography.md} & D1 7,549 (\val{d1} with the 26A434 decoder), D2 \val{d2}; ten tensor MMA opcodes admitted & \cref{sec:opcode-identity-counts,sec:admission-decodes-executes} \\
The project framer's first check walked 112,952
instructions where Apple's decoder finds 85,806 & \code{docs/agx3-oracle.md} section 10 & the framer agrees with Apple on 1,725 of 22,187 constant programs and
misframes 200 of 200 current programs; Apple's decoder
is the ground truth & \cref{sec:framing}, MM~\mm{25.88}, MM~\mm{25.209} \\
Encoder held-out re-authoring 28.2 percent (before seven fixes) & MM~\mm{25.88} & 80.8 percent & \cref{sec:field-maps}, MM~\mm{25.88} \\
\Op{12674} fill-slot field is operand 1 bits
20–27 & MM~\mm{6} & bits 20–23 fill slot (slot + 1), bits 24–31 the load's
own wait mask & \cref{sec:scoreboard-publish-wait} \\
Constant-program binding reads at "halfwords 8 + K and 10 + K" (K the
halfword offset, MM~\mm{25.179}) against "8 + 4k and 10 + 4k for the
k-th buffer" (MM~\mm{25.59}) & MM~\mm{25.179}, MM~\mm{25.59} & the same rule with K = 4k & \cref{sec:constant-program-record} \\
Constant-program metadata read as a register-spill description & MM~\mm{25.113} & it describes the constant program & \cref{sec:constant-program-record}, MM~\mm{25.120} \\
The "76 inert" list itemised 92 bytes & MM~\mm{25.148} (ledger) & 76 inert bytes; the descriptor tail (392, 16) is separate and not one of
the 148 class constants & \cref{sec:object-library-archive}, MM~\mm{25.148} \\
"The image is ours" census covering one novel kernel & MM~\mm{25.148} (ledger) & the 421-archive delivered-image census, every byte assigned & \cref{sec:object-library-archive}, MM~\mm{25.148} \\
\bottomrule
\end{xltabular}

\subsection*{Scalar arithmetic, control flow, special registers, memory, cross-lane
and synchronization}\label{sub:scalar-arithmetic-control}

\begin{xltabular}{\linewidth}{@{}>{\hsize=1.10\hsize}X>{\hsize=0.63\hsize}X>{\hsize=1.44\hsize}X>{\hsize=0.83\hsize}X@{}}
\caption{Superseded counts and wordings about scalar arithmetic, control flow, special registers, memory, cross-lane operations and synchronization.}\label{tab:acct-scalar-memory}\\
\toprule
Old reading or accounting & Where & Current reading & Where \\
\midrule
\endfirsthead
\tablecontinued{4}
\toprule
Old reading or accounting & Where & Current reading & Where \\
\midrule
\endhead
"fsat" names \op{1062}; "saturate returns
0" & early map; the anomaly came from a record that ran one lane & \opn{1062} is a cross-lane \code{dfdx} with saturation (function
unconfirmed); saturate is \op{904}, sat(a + K) & MM~\mm{25.90}; \cref{sec:fp32-arithmetic,sec:other-cross-lane-forms} \\
Printed integer condition code 15 is "something other than lt/gt/signed" & \code{isa/g17-condition-codes.toml} notes & it is the bit test (a \& b) != 0, field code 1 of \op{11375} & MM~\mm{25.90}; \cref{sec:compares-flags-selects} \\
\code{(uint)x} lowered as a 34-instruction decomposition that decoded
the float's bits; \code{(uint)(-5.5)} gave 5 & this compiler, before MM~\mm{25.186} & \op{9320}, truncate and saturate, gives 0;
Metal leaves the case undefined and each emulation matched its own
program's GPU output (a compiler difference, not a
hardware error) & MM~\mm{25.186}; \cref{sec:conversions} \\
Integer-to-float conversion code 4 for every source (inherited from an
unsigned witness; a negative int32 became a value near $2^{32}$) & \code{agxforge/g17/cc.py} before MM~\mm{25.90} & operand 2 is signedness: 4 unsigned, 5 signed (2/3 on the 16-bit form) & MM~\mm{25.90}, MM~\mm{25.195}; \cref{sec:conversions} \\
The join was called \code{exec.restore} & early \code{isa/g17-scalar-isa.toml} & \op{577} with count 1 & \code{agxforge/g17/cf.py}; \cref{sec:execution-mask-family} \\
A 4-byte backward branch \code{de 6f f3 1e} with a 12-bit displacement
whose sign is word bit 28 & \code{isa/g17-scalar-isa.toml} & the first four bytes of the 10-byte \op{458}: kind 2 with b2.4 and b2.0 set; d1-d10 are the 10-byte
map's weights 2 to $2^{10}$ in bytes 0–2; the "sign" b3.4
is the $2^{14}$ weight, reading as a sign only because a small negative
displacement sign-extends through it & \code{agxforge/g17/cf.py}; \cref{sec:execution-mask-family} \\
Byte 0x82 of the 4-byte special-register read is
\code{quadgroup\_index\_in\_threadgroup}; 0x98 is
\code{simdgroups\_per\_threadgroup} (from kernels said to differ only in
that byte) & early \code{isa/g17-scalar-isa.toml} & the toml recorded registers the builtins' kernels read;
0x82 is the lane index (\code{SR\_SIMD\_ELEM}), 0x98 threads-per-threadgroup x & MM~\mm{0.8}, \code{isa/g17-special-register-numbers.txt},
\code{agxforge/g17/mdgen.py}; \cref{sec:special-registers} \\
\code{SR\_SIMD\_GRP} "not generalized", no slot-29 entry & MM~\mm{0.8} table & entry 53 is exactly the SR133 read added by the multi-simdgroup compile
(\code{execution\_simdgroups} sweep); single-builtin compiles carry
{[}53{]}; entry 152 → 4 added later & MM~\mm{25.93}, MM~\mm{25.149.1}; 
\cref{sec:slot-29-metadata-numbering} \\
A store reading a special-register value directly stores what has not
arrived (ledger, 2026–09-08); every read is a late value needing a wait & ledger; MM~\mm{6}, MM~\mm{0.8}, MM~\mm{25.90} & the failing stores named one register as value and index with the index
slot releasing it; a slot-0 read needs no wait (latency probe
2026–09-24, 90 of 90 rows reproduced) & MM~\mm{25.117}, MM~\mm{25.121}; \cref{sec:read-latency-whether} \\
The BIF0 read page fault came from the tensor-probe configuration the
rule names & MM~\mm{17} rule 34 before its 2026–09-23 correction & that configuration was a different process (pid 76676), which hung the
GPU for about 320 s and is known to hang, not to fault; the fault came
from a \code{regprobe\_multi} kernel whose configuration was never
recorded (corrected from another lane's reading of the
crash report and system log) & MM~\mm{17} rule 34; \cref{sec:out-of-range-accesses-observed} \\
The 8-byte \op{17229} "has no byte9 and no wait
bit" & \code{agxforge/g17/cc.py} & the encoding carries wait bits: 15 corpus instances match their pending
slots, and Apple's \code{tex2d-read} carries 0x1000010
on one; the statement describes this compiler's encoder
only & \code{isa/g17-wait-mask-census.json}, MM~\mm{25.146}; 
\cref{sec:scoreboard} \\
"Eight raw operation codes of \op{10090} unassigned by any Metal intrinsic" & \code{isa/g17-atomic-semantics.json} & read off one witness's field placement; retracted; the
fourth operation bit of the 12-byte form is byte11{[}4{]}, and
unassigned codes must be counted per form & \code{isa/g17-atomic-family.toml}; \cref{sec:atomics} \\
\bottomrule
\end{xltabular}

\subsection*{The tensor unit, performance and submission}\label{sub:tensor-unit-performance}

\begin{xltabular}{\linewidth}{@{}>{\hsize=1.31\hsize}X>{\hsize=0.59\hsize}X>{\hsize=1.33\hsize}X>{\hsize=0.77\hsize}X@{}}
\caption{Superseded counts and wordings about the tensor unit, performance and submission.}\label{tab:acct-tensor-submission}\\
\toprule
Old reading or accounting & Where & Current reading & Where \\
\midrule
\endfirsthead
\tablecontinued{4}
\toprule
Old reading or accounting & Where & Current reading & Where \\
\midrule
\endhead
Two declaration counts for the tensor family: 132 opcodes in scheduling
classes 170–175, against 535 declarations in opcode ranges 5090–5618 and
10384–10389 (18 forms with an 8-word destination, 5090–5107, 8 admitted;
36 with a 4-word destination, 5108–5143; 6 integer forms, 10384–10389, 2
admitted; 475 \code{\_shifted} forms, 5144–5618) & MM~\mm{2}, MM~\mm{0.3}; MM~\mm{25.6}, MM~\mm{25.4/F4} "F4, narrowed" & The two counts use different bases and do not conflict. The 132 counts by
scheduling class; the 535 counts by opcode range and includes every \code{\_shifted} declaration.
Ten forms are decoder-admitted in either accounting & \cref{sec:instruction-family} \\
"None of the 56 MMA forms the decoder tables declare" has a scale
operand & MM~\mm{11} & The statement stands. How the 56 relates to the 60 non-\code{\_shifted} declarations
(18 + 36 + 6) is not stated in the sources & \cref{sec:instruction-family} \\
Destination keep-flag rule quoted over three populations: 2,192 MMAs in
140 programs (1,093 flagged); 540,396 flag-set instances with 155,963
converse counterexamples; 758,739 "instances" & MM~\mm{25.22}; MM~\mm{6}; MM~\mm{0.4} & All three agree on the one-way rule (set implies the next MMA reuses the
tuple) with zero counterexamples. The 758,739 population is not defined
by its source and is not 540,396 + 155,963 (696,359) & \cref{sec:modifiers-keep-flag} \\
"Whether release clears or frees storage is open"; "a released register
reads 0 afterwards" & MM~\mm{25.89}; earlier text of \cref{sec:modifiers-keep-flag} & A release acts per lane under the execution mask; a released register
read zero in every controlled run, but what it reads is not established
in general; nothing reclaimed it within the measured bound & \cref{sec:release-does} (MM~\mm{25.95}, MM~\mm{25.145}) \\
\op{5100} and \op{5101} "are never
compiler-emitted" (from recon sections 27 and 104) & MM~\mm{19} & Emitted under \code{relaxed\_precision} & \cref{sec:it-reached} (MM~\mm{17} rule 29) \\
Element code 4 does not mean a 32-bit access (the withdrawal of "code 4
means 32-bit") & MM~\mm{25.5/R7} "R7 narrowed" & The element code is the access width in bytes; code 4 is emitted by fp32
accesses (two witnesses) & \cref{sec:tensor-memory-path} (MM~\mm{25.37}) \\
\op{435} and \op{437} read as
shifts & earlier opcode notes & They are ANDs (\opn{435} executed over 4,096 lanes, MM~\mm{25.192}) & \cref{sec:quantized-formats} (MM~\mm{24.5}) \\
Apple \code{matmul2d} at M512 took 58.6~µs (round one); this compiler
1.17× behind, unrolled 2.1–4.5× & MM~\mm{25.98} & Round one did not reproduce; interleaved medians give 106.6~µs for Apple
at M512, the K loop 1.07× behind & \cref{sec:measured-bottlenecks-studied} (MM~\mm{25.107}) \\
Residency of a used-125 program: about 43.6 per core (v3), then about 58
per core (v4), and a shortfall at S = 1,024 (v5) & MM~\mm{25.115} & v3 ran beside an orphaned kernel; the v5 shortfall was a cold-clock
effect. Warm, 64–69 per core from 49 to 125 registers, 95.0 at 19 & \cref{sec:residency-occupancy-measurements} (MM~\mm{25.121}) \\
Instruction-fetch step as a whole-kernel size bracket (16,066, 17,158{]}
bytes & recon section 156; MM~\mm{19} & The variable is the loop body: 15,971 bytes fast, 17,247 slow,
bracketing 16 KiB & \cref{sec:memory-bandwidth-footprint} (MM~\mm{19}) \\
GPU performance counters "not reachable without Apple's
GUI or entitlement" (the "M4 status" paragraph) & MM~\mm{25.144.4} & 89 counters are recorded by \code{xctrace} with the GUI-saved
"Performance Limiters" template, later in the same section;
Metal's public API still offers timestamps only & \cref{sec:dispatch-encoder-command-buffer} \\
"Power was not measured" & MM~\mm{20} "Hardware identity" & Measured on the GPU rail; MM~\mm{20}'s own later
paragraph says so & \cref{sec:clocks-power-gpu} (MM~\mm{25.56}, MM~\mm{25.60}) \\
Below Metal, "three command buffers make one selector-7 call", a
persistent command stream with no MMIO doorbell & MM~\mm{0.15}; earlier text of 
\cref{sec:dispatch-encoder-command-buffer} & Selectors 7, 16 and 28 are called in queue creation; each Submit enters
the kernel through \code{IOConnectTrap4} (one record) or
\code{IOConnectCallMethod} selector 0x1d (several records) & \cref{sec:opening-device-contexts,sec:submit-record-call} \\
Submit kicks the GPU only through shared-memory writes, with no
per-dispatch kernel call; a candidate store at Submit+0x11c & MM~\mm{25.143}; README row 6a & Each Submit enters the kernel through the trap or the method call; the
store at Submit+0x11c is an output-argument write & \cref{sec:submit-record-call} \\
The \code{FinalizeShmemHeader} ready bit is the doorbell & README row 3 & Disproved by a live control; see the Submit path above & \cref{sec:submit-record-call} \\
\bottomrule
\end{xltabular}
